\chapter{Summary and further scope}\label{Summary}

Simulations from the 3-D Hele--Shaw model present a comprehensive view of the evolution of systems with deformable particles through hydrodynamic interactions in confined parallel-wall geometries. First, it remains consitent with the 2-d model regarding alignment of paticles along the flow-direction via the Hele--Shaw quadrupolar interactions\cite{Reference33}. The tendency of particles to align in the flow-direction was one of the first key observations reported by all the experimental studies\cite{Reference15}\cite{Reference16}\cite{Reference30} reported in Chapter.~\ref{experimental_observations}. Second, the 3-D simulations exhibited the tendency of particles to remain stationary, relative to its neighbours, at fixed interparticle separations after reaching stagnation. This was another important feature observed in the experiments. Lastly, the 3-D simulations enrich the understanding of the particle dynamics with its ability to record the hydrodynamic interactions between particles and their migration in the flow-graident plane. Evolutionary features exhibited for a four-particle system had strong resemblances to that of the direct-numerical-simulations both in the flow and gradient directions thereby reinforcing the fact that the hydrodynamic interactions responsible for the key features of particle evolution in experiments were identified and that the simple Hele--Shaw model was able to predict the particle trajectories in a multiparticle system.

In future, with the ability to track motion in the flow-vorticity and flow-gradeint plane, using this model one can further the understanding and recreation of important features displayed by bigger systems over a range of volume-fractions in stokes regime both for mono-disperse and poly-disperse systems. Efforts can be made to extend the model to reflect the particle evolution in flows with inertia, where transverse motion of particles was quite prevalent, and also in the case of non-newtonian systems.